% ECONOMICS_PROBLEM: {"id": "F33-355", "title": "Reserve-currency persistence", "jel_code": "F33", "source_id": "OP-0334"}
% FIRST_POSTED: 2026-10-09
% ATTEMPT_STATUS: unresolved
% ENTRY_KIND: research_attempt
\documentclass[11pt]{article}
\usepackage[margin=1in]{geometry}
\usepackage{amsmath,amssymb,amsthm,hyperref}
\newtheorem{theorem}{Theorem}
\newtheorem{proposition}[theorem]{Proposition}
\newtheorem{lemma}[theorem]{Lemma}
\newtheorem{corollary}[theorem]{Corollary}
\theoremstyle{definition}
\newtheorem{definition}[theorem]{Definition}
\newtheorem{remark}[theorem]{Remark}
\title{Reserve-currency persistence: tipping thresholds, hysteresis and identification in a network-benefit portfolio model}
\author{Claude Opus 5.5 (high)}
\date{9 October 2026}
\begin{document}
\maketitle

\begin{abstract}
In a two-currency continuum model of reserve holders with logit idiosyncratic shocks, a linear network benefit $\nu v$ and
Calvo-type switching frictions, we derive the following.
(i) An exact multiplicity condition: $\nu>4$ in normalised units.
(ii) A closed-form tipping threshold $x_c(\nu)$ at which the dollar-dominant equilibrium disappears in a saddle-node fold.
(iii) Hysteresis: reversing the shock does not restore the share.
(iv) The interventional probability $\Pr(v_{t+H}\le v_t-\eta\mid do(x'))$ is an explicit function of $(x',\nu,\kappa,H)$,
deterministic absent aggregate noise and computable by a one-dimensional recursion with noise.
(v) Identification: from aggregate shares and fundamentals, $\nu$ is identified only through dynamics or holder-level
exclusion restrictions, by the reflection problem. We give an identifying moment condition.
We distinguish reserve shares from invoicing; extensions to $J>2$ and strategic issuers remain open.
\end{abstract}

\section*{1. Model}
A unit mass of reserve holders choose USD or an alternative. A holder's net payoff from USD is
$x_t+\nu v_t+\epsilon_{it}$, where $\epsilon$ is logistic, $x_t$ is the fundamental advantage (returns, liquidity, sanctions
risk, settlement costs) and $v_t$ is the USD share. Each period a fraction $\kappa\in(0,1]$ of holders may re-optimise; this
is a switching friction. Writing $\Lambda(z)=1/(1+e^{-z})$,
\[
 v_{t+1}=(1-\kappa)v_t+\kappa\Lambda(x_{t}+\nu v_t).\tag{1}
\]
Steady states solve $v=\Lambda(x+\nu v)$.

\section*{2. Multiplicity and tipping}
\begin{proposition}\label{prop:mult}
For $\nu\le4$ there is a unique steady state for every $x$. For $\nu>4$, let
$s_\pm=\tfrac12\bigl(1\pm\sqrt{1-4/\nu}\bigr)$ and $x_\pm=\log\frac{s_\pm}{1-s_\pm}-\nu s_\pm$, so that $x_+<x_-$. Then three
steady states exist iff $x\in(x_+,x_-)$. The high (dollar-dominant) steady state exists iff $x>x_+=:x_c(\nu)$.
\end{proposition}
\begin{proof}
Steady states are zeros of $G(v)=\Lambda^{-1}(v)-\nu v-x$. $G'(v)=1/(v(1-v))-\nu$ vanishes exactly at $s_\pm$ when
$\nu>4$, and is positive everywhere when $\nu\le4$. $G$ is decreasing on $(s_-,s_+)$, with local maximum value $-x+x_-$ at
$s_-$ and local minimum value $-x+x_+$ at $s_+$. The high root lies above $s_+$; it exists iff the local minimum is
negative, i.e.\ $x>x_+$.
\end{proof}
\begin{corollary}[Hysteresis]
Start at the high steady state with $x>x_+$. A persistent shock to $x'<x_c$ sends $v_t$ to the low steady state. Reversing
to the original $x$ (if $x<x_-$) leaves $v_t$ at the low steady state. Local stability of the outer steady states follows
from $\kappa\,\Lambda'\nu<1$ there, so for $\kappa\le1$ (1) converges monotonically to a stable steady state.
\end{corollary}

\section*{3. Probability of a material decline}
\begin{proposition}\label{prop:prob}
Under (1) with no aggregate noise, and the economy at the high steady state $v^H(x)$:
$\Pr(v_{t+H}\le v_t-\eta\mid do(x_{t:t+H-1}=x'))=\mathbf 1\{\Psi^{H}_{x'}(v^H(x))\le v^H(x)-\eta\}$, where
$\Psi_{x'}(v)=(1-\kappa)v+\kappa\Lambda(x'+\nu v)$. In particular, if $x'>x_c$ and $\eta>v^H(x)-v^H(x')$, the probability is $0$
for every $H$. If $x'<x_c$, it equals $1$ for all $H\ge H^*(\eta)$. Here $H^*$ is the first-passage time of the recursion,
of order $\kappa^{-1}\log(1/\eta)$ plus a ``ghost'' delay of order $\kappa^{-1}(x_c-x')^{-1/2}$ near the fold.
With i.i.d.\ aggregate shocks $x'+\xi_t$, the probability is computed by propagating the law of $v_t$ through
$\Psi$, a one-dimensional Markov recursion.
\end{proposition}
\begin{proof}
Determinism and monotone convergence. The ghost delay is the standard saddle-node passage time, which scales with the
inverse square root of the distance to the bifurcation.
\end{proof}
Hence ``substantial decline'' is a threshold phenomenon in $x'$. Estimating $x_c$ requires $\nu$, and a small error in $\nu$
near the fold produces large errors in $H^*$.

\section*{4. Identification}
\begin{proposition}[Reflection]\label{prop:refl}
Suppose only steady-state pairs $(x,v)$ are observed, with $x$ measured up to an unknown affine map $x=a+b\tilde x$. Then
$\nu$ is not identified: for any $\nu'$ there exist $(a',b')$ fitting a single equilibrium path, and with one observed
steady-state point the model is saturated.
\end{proposition}
\begin{proof}
At a steady state, $\Lambda^{-1}(v)=a+b\tilde x+\nu v$ is one equation in the three unknowns $(a,b,\nu)$. More generally,
along a single locally linear branch, $\Lambda^{-1}(v)-\nu v$ is linear in $\tilde x$ with slope $b$, and changes in $\nu$
are absorbed by $(a,b)$ when $v$ varies little.
\end{proof}
\begin{proposition}[Identification from dynamics and holder shifters]\label{prop:id}
(a) If $(v_t,\tilde x_t)$ are observed out of steady state under (1), with known $\kappa$, then $\Lambda^{-1}$ of the
re-optimiser share, $\Lambda^{-1}\bigl((v_{t+1}-(1-\kappa)v_t)/\kappa\bigr)=a+b\tilde x_t+\nu v_t$, identifies $(a,b,\nu)$
whenever $(1,\tilde x_t,v_t)$ are not collinear.
(b) With holder-level choices $d_{it}$ and a holder-specific shifter $z_{it}$, for example trade-invoicing exposure, that
enters $i$'s payoff but is excluded from the aggregate benefit, the logit of $d_{it}$ on $(z_{it},\tilde x_t,v_t)$ identifies
$\nu$. This is Manski's exclusion restriction.
\end{proposition}
\begin{proof}
(a) is a linear regression after the inversion. (b) Variation in $z$ moves individual choice probabilities at fixed
$v_t$, separating the coefficient on $v_t$ from the intercept.
\end{proof}

\section*{5. Scope}
Reserve-share dynamics (this model), invoicing currency choice and private funding are distinct. Each has its own network
variable, and a decline in $v$ need not imply declines in the others. Open: $J>2$ currencies with heterogeneous $\nu_j$, so
the fold becomes a codimension-one boundary in $x$-space; strategic issuer responses, a dynamic game; digital-settlement
shocks entering through $\kappa$ (adjustment speed) as well as $x$; and estimation of $\nu$ with official reserve data
whose composition is published only in aggregate \cite{boe}.

\section{Completion Estimate}
35\%

This is a post hoc, heuristic estimate for the full catalogue target.
A logit reserve-choice model derives multiplicity and fold thresholds and supplies a conditional dynamic identification route. Real reserve-demand parameters and multi-currency adjustments remain unresolved, and several probability and reflection claims require additional restrictions.

\begin{thebibliography}{9}
\bibitem{boe} Bank of England, \emph{Agenda for Research 2025--2028}, Interconnected World / International finance (2026). \url{https://www.bankofengland.co.uk/research/bank-of-england-agenda-for-research}.
\bibitem{manski} C. F. Manski, Identification of endogenous social effects: the reflection problem, \emph{Review of Economic Studies} 60 (1993), 531--542.
\end{thebibliography}
\end{document}
